\documentclass[10pt]{elegantnote}
\usepackage{graphicx}
\usepackage{unicode-math}
\usepackage{framed} %段落加边框
\usepackage{makecell} %表格内强制换行
\usepackage{titletoc} % 章节目录
\usepackage{tabularx}
\usepackage{siunitx} % 数字，单位宏包
\sisetup{
    separate-uncertainty = true,
    inter-unit-product = \ensuremath{{}\cdot{}}
  } %单位间加小圆点
\usepackage{array}

\setmathfont{XITS Math}

% 取消pdf的压缩，最终版取消这个命令即可
\special{dvipdfmx:config z 0}

%自定义命令 颜色
\definecolor{epubblue}{RGB}{1,126,218}
\newcommand*{\cepubblue}{\textcolor{epubblue}}

% 字体

\RequirePackage[fontset=none]{ctex}
\setmainfont{Times New Roman}
\setCJKmainfont{SimSun}[BoldFont=FandolSong Bold, AutoFakeSlant]
\setCJKsansfont[BoldFont={FZHei-B01}]{FZKai-Z03}
\setCJKmonofont[BoldFont={FZHei-B01}]{FZFangSong-Z02}
\setCJKfamilyfont{zhsong}{FZShuSong-Z01}
\setCJKfamilyfont{zhhei}{FZHei-B01}
\setCJKfamilyfont{zhkai}[AutoFakeBold]{FZKai-Z03}
\setCJKfamilyfont{zhfs}[BoldFont={FZHei-B01}]{FZFangSong-Z02}
\setCJKfamilyfont{yahei}{Microsoft YaHei}
\setmathfont{XITS Math}

% 代码
\RequirePackage{listings}
%Python代码排版设置
\lstset{
    language=Python, % 设置语言
 breaklines=true, % 自动换行
 backgroundcolor=\color[rgb]{0.96,0.96,0.96},% 背景颜色
 keywordstyle=\bfseries\color{RoyalBlue}, % 设置关键字为粗体，颜色为 RoyalBlue
 morekeywords={syms,eqns,solve,simplify}, % 设置更多的关键字，用逗号分隔
 emph={left,right,mid}, % 指定强调词，如果有多个，用逗号隔开
    emphstyle=\bfseries\color{Rhodamine}, % 强调词样式设置
    commentstyle=\color{black!50!white}, % 设置注释样式，斜体，浅灰色
    stringstyle=\color{PineGreen!90!black}, % 设置字符串样式
    columns=flexible,
    numbers=left, % 显示行号在左边
    numbersep=2em, % 设置行号的具体位置
    numberstyle=\footnotesize, % 缩小行号
    xleftmargin=1em, % 左边距
    xrightmargin=1em, % 右边距
    aboveskip=1em, % 上边距
    % framexleftmargin=2em, % 行号在框内
    framesep=1em, % 设置边框与代码的距离
    frame=shadowbox,rulesep=3pt,rulesepcolor=\color{red!20!green!20!blue!20} % 阴影边框
    }


\setlist[itemize,1]{\textcolor{ecolor}{\label =\bullet} }
\setlist[itemize,2]{\textcolor{ecolor}{\label =$\circ$} }

\begin{document}

\begin{sloppypar} %文字溢出自动换行


\begin{center}
    \bfseries \Large \textcolor{ecolor} {《计算思维与数学建模》}

    \bfseries \Large \textcolor{ecolor} {教学设计}

    \bfseries \Large \textcolor{ecolor} {张斌}
\end{center}

% \tableofcontents

% \newpage

\section{教学目标}
    \begin{enumerate}
        \item 知识目标
            \begin{itemize}
                \item 了解 MATLAB 的基本用法，包括常用的命令，脚本结构，程序的分节与注释，矩阵的基本操作，绘图方式；
                \item 熟悉 MATLAB 的程序结构，编程模式，能自由地创建和引用函数；
                \item 熟悉常见模型的求解算法和套路。
            \end{itemize}
        \item 能力目标
            \begin{itemize}
              \item 能够熟练运用 MALTAB 软件，能选择合适的方法建立数学模型和求解。
            \end{itemize}
        \item 素质目标
            \begin{itemize}
              \item 具有一定的表达能力，能清晰的描述所建立的模型及其结果；具备规范的撰写数学建模论文的能力；具有较好的沟通能力，掌握团队协作技能，保证成员共同完成任务目标。
            \end{itemize}
\end{enumerate}

\begin{table}[htpb]
    \centering
    \caption{课程目标对专业毕业要求的支撑}
    \begin{tabular}{cccc}
        \bottomrule
        毕业目标 & 课程目标 & 支撑度 &  \\ \hline
        \begin{tabular}[c]{@{}c@{}}毕业目标1-1:掌握必要的数学和物理知识，\\ 能够针对复杂化工问题进行建模和求解。 \end{tabular} & \begin{tabular}[c]{@{}c@{}}目标1\\ 目标2\\ 目标3 \end{tabular} & M &  \\ \toprule
    \end{tabular}
\end{table}

\section{教学安排}

本课程共1.5学分，48学时实践，分散3周进行。

% \begin{table}[htpb]
%     \centering
%     \begin{tabular}{cccc}
%         \bottomrule
%         时间 & 章节 & 教学内容 & 教学形式 \\ \hline
%         第一周 & Matlab基础 & 界面功能、数据类型、符号运算 & 线下教学 \\ \hline
%         第二周 & 程序设计方法 & 程序控制结构、关系、逻辑运算 & 线下教学 \\ \hline
%         第三周 & 数学建模基础 & 分析抽象、建模步骤、初等模型 & 线下教学 \\ \hline
%         第四周 & 常用建模函数 & 插值函数，曲线拟合，参数传递 & 线下教学 \\ \hline
%         第五周 & 简单优化模型 & 优化目标、优化策略、函数极值 & 线下教学 \\ \hline
%         第六周 & 数学规划模型 & 线性规划模型、非线性规划模型 & 线下教学 \\ \hline
%         第七周 & 微分方程模型 & 导数、极限、定积分、不定积分 & 线下教学 \\ \hline
%         第八周 & 智能优化算法 & 遗传算法，免疫算法，动态规划 & 线下教学 \\ \toprule
%     \end{tabular}
% \end{table}

% \begin{table}[htpb]
%     \centering
%     \begin{tabular}{@{}|c|c|c|c|c|@{}}
%     \bottomrule
%         & 学时 & 线下教学(32)                                                                                                             & 学时 & 线上学习(16)                                                       \\ \hline
%         \begin{tabular}[c]{@{}c@{}}第1周\\ Matlab基础\end{tabular} & 4  & \begin{tabular}[c]{@{}c@{}}Matlab的界面功能\\      Matlab中的数据类型\\      Matlab矩阵的操作\\      Matlab中的运算\end{tabular}         &    &                                                                \\ \hline
%         \begin{tabular}[c]{@{}c@{}}第1周\\ Matlab程序设计\end{tabular} & 4  & \begin{tabular}[c]{@{}c@{}}结构化程序设计方法\\      流程控制语句\\      关系、逻辑运算\\      函数文件与脚本文件\end{tabular}                      &    &                                                                \\ \hline
%         \begin{tabular}[c]{@{}c@{}}第2周\\ 数学建模基础\end{tabular} & 4  & \begin{tabular}[c]{@{}c@{}}对实际问题的分析、抽象过程\\      数学建模的基本步骤\\      用初等方法建立模型\\      双层玻璃窗保温原理的数学模型\end{tabular}        & 4  & \begin{tabular}[c]{@{}c@{}}划船比赛的成绩\\      节水洗衣机\end{tabular}   \\ \hline
%         \begin{tabular}[c]{@{}c@{}}第2周\\ 常用建模函数\end{tabular} & 4  & \begin{tabular}[c]{@{}c@{}}一维插值，二维插值，样条插值\\      多项式拟合\\      非线性曲线拟合\\      参数传递函数\end{tabular}                     &    &                                                                \\ \hline
%         \begin{tabular}[c]{@{}c@{}}第2周\\ 简单优化模型\end{tabular} & 4  & \begin{tabular}[c]{@{}c@{}}优化目标\\      优化策略\\      函数的极值\\      易拉罐形状和尺寸的最优设计\end{tabular}                           & 4  & \begin{tabular}[c]{@{}c@{}}生猪的出售时机\\      存储模型\end{tabular}    \\ \hline
%         \begin{tabular}[c]{@{}c@{}}第3周\\ 数学规划模型\end{tabular} & 4  & \begin{tabular}[c]{@{}c@{}}数学规划模型的概念\\      线性规划的求解方法\\      非线性规划模型\\      自来水运输问题\end{tabular}                     & 4  & \begin{tabular}[c]{@{}c@{}}奶制品的生产计划\\      汽车生产计划\end{tabular} \\ \hline
%         \begin{tabular}[c]{@{}c@{}}第3周\\ 微分方程模型\end{tabular} & 4  & \begin{tabular}[c]{@{}c@{}}Matlab求函数的导数、极限\\      Matlab求定积分、不定积分\\      Matlab求解微分方程\\      药物在体内的分布模型\end{tabular} & 4  & \begin{tabular}[c]{@{}c@{}}人口增长模型\\      传染病模型\end{tabular}    \\ \hline
%         \begin{tabular}[c]{@{}c@{}}第3周\\ 智能优化算法\end{tabular} & 4  & \begin{tabular}[c]{@{}c@{}}遗传算法\\ 免疫算法\\ 动态规划\\ 最短路径规划问题\\ 数学建模论文撰写规范\end{tabular}                                   &    &                                                                \\ \toprule
%     \end{tabular}
%     \end{table}

\begin{table}[htpb]
    \centering
    \begin{tabular}{@{}|c|c|c|c|c|@{}}
    \bottomrule
        & 学时 & 线下教学(32)                                                                                                             & 学时 & 线上学习(16)                                                       \\ \hline
        \begin{tabular}[c]{@{}c@{}}第1周\\ Matlab基础\end{tabular} & 4  & \begin{tabular}[c]{@{}c@{}}Matlab的界面功能\\      Matlab中的数据类型\\      Matlab矩阵的操作\\      Matlab中的运算\end{tabular}         &  2  &  \begin{tabular}[c]{@{}c@{}}结构化程序设计方法\\      流程控制语句\\      关系、逻辑运算\\      函数文件与脚本文件\end{tabular}                               \\ \hline
        
        \begin{tabular}[c]{@{}c@{}}第1周\\ 初等模型\end{tabular} & 4  & \begin{tabular}[c]{@{}c@{}}      数学建模的基本步骤\\  常用建模函数\\    用初等方法建立模型\\      双层玻璃窗保温原理数学模型\end{tabular}        & 2  & \begin{tabular}[c]{@{}c@{}}划船比赛的成绩\\      节水洗衣机\end{tabular}   \\ \hline
        
        \begin{tabular}[c]{@{}c@{}}第2周\\ 简单优化模型\end{tabular} & 4  & \begin{tabular}[c]{@{}c@{}}优化目标\\      优化策略\\      函数的极值\\      易拉罐形状和尺寸的最优设计\end{tabular}                           & 2  & \begin{tabular}[c]{@{}c@{}}生猪的出售时机\\      存储模型\end{tabular}    \\ \hline
        \begin{tabular}[c]{@{}c@{}}第2周\\ 数学规划模型\end{tabular} & 4  & \begin{tabular}[c]{@{}c@{}}数学规划模型的概念\\      线性规划的求解方法\\      非线性规划模型\\      自来水运输问题\end{tabular}                     & 2  & \begin{tabular}[c]{@{}c@{}}奶制品的生产计划\\      汽车生产计划\end{tabular} \\ \hline
        \begin{tabular}[c]{@{}c@{}}第2周\\ 微分方程模型\end{tabular} & 4  & \begin{tabular}[c]{@{}c@{}}Matlab求函数的导数、极限\\      Matlab求定积分、不定积分\\      Matlab求解微分方程\\      药物在体内的分布模型\end{tabular} & 2  & \begin{tabular}[c]{@{}c@{}}人口增长模型\\      传染病模型\end{tabular}    \\ \hline
        \begin{tabular}[c]{@{}c@{}}第3周\\ 离散模型\end{tabular} & 4  & \begin{tabular}[c]{@{}c@{}}多属性决策问题\\      层次分析法\\      汽车选购\end{tabular} & 2  & \begin{tabular}[c]{@{}c@{}}席位分配\end{tabular}    \\ \hline
        \begin{tabular}[c]{@{}c@{}}第3周\\ 统计模型\end{tabular} & 4  & \begin{tabular}[c]{@{}c@{}}线性回归\\      非线性回归\\      自回归\\      logit回归\end{tabular} & 2  & \begin{tabular}[c]{@{}c@{}}判别分析模型\\主成分分析模型\end{tabular}    \\ \hline
        \begin{tabular}[c]{@{}c@{}}第3周\\ 智能优化算法\end{tabular} & 4  & \begin{tabular}[c]{@{}c@{}}遗传算法\\ 免疫算法\\ 神经网络 \\ 数学建模论文撰写规范\end{tabular}                                   &  2  &   \begin{tabular}[c]{@{}c@{}} 动态规划\\ 最短路径规划问题\end{tabular}                                                      \\ \toprule
    \end{tabular}
    \end{table}

\newpage
\section{考核评价}

学生的成绩由个人成绩(40\%)和团队成绩(60\%)组成。

个人考核形式为完成测试。

团队考核形式为以小组为单位完成分组任务。

\begin{table}[htpb]
    \centering
    \caption{考核评价方式及权重}
    \begin{tabular}{cccc}
        \bottomrule
        & 个人考核(40\%) & 团队考核(60\%) &  \\ \hline
        Matlab基础 & 自测题(5\%) & 利用Matlab求解反应平衡常数(5\%) &  \\ \hline
        初等模型 & 自测题(5\%) & 三层玻璃窗保温原理的数学模型(5\%) &  \\ \hline
        简单优化模型 & 自测题(5\%) & 易拉罐球台模型的最优尺寸设计(5\%) &  \\ \hline
        数学规划模型 & 自测题(5\%) & 有限制的自来水运输模型规划(5\%) &  \\ \hline
        微分方程模型 & 自测题(5\%) & 一室模型多次重复给药的方案(5\%) &  \\ \hline
        离散模型 & 自测题(5\%) & 奖学金申请(5\%) &  \\ \hline
        统计模型 & 自测题(5\%) & 蛋白质含量对婴儿身高影响(5\%) &  \\ \hline
        智能优化算法 & 自测题(5\%) & 完成一篇完整的数学建模论文(25\%) &  \\ \toprule
    \end{tabular}
\end{table}


% \section{教学大纲}

% \begin{center}
%     第一周\ Matlab基础
% \end{center}

% 主要内容：

% \begin{enumerate}
%   \item Matlab的界面功能
%   \item Matlab中的数据类型
%   \item Matlab矩阵的操作
%   \item Matlab中的运算
% \end{enumerate}

% \begin{center}
%     第二周\ 程序设计方法
% \end{center}

% 主要内容：

% \begin{enumerate}
%   \item 结构化程序设计方法
%   \item 流程控制语句
%   \item 关系、逻辑运算
%   \item 函数文件与脚本文件
% \end{enumerate}

% \begin{center}
%     第三周\ 数学建模基础
% \end{center}

% 主要内容：

% \begin{enumerate}
%   \item 对实际问题的分析、抽象过程
%   \item 数学建模的基本步骤
%   \item 用初等方法建立模型
%   \item 双层玻璃窗保温原理的数学模型
% \end{enumerate}

% \begin{center}
%     第四周\ 常用建模函数
% \end{center}

% 主要内容：

% \begin{enumerate}
%   \item 一维插值，二维插值，样条插值
%   \item 多项式拟合
%   \item 非线性曲线拟合
%   \item 参数传递函数
% \end{enumerate}

% \begin{center}
%     第五周\ 简单优化模型
% \end{center}

% 主要内容：

% \begin{enumerate}
%   \item 优化目标
%   \item 优化策略
%   \item 函数的极值
%   \item 易拉罐形状和尺寸的最优设计
% \end{enumerate}


% \begin{center}
%     第六周\ 数学规划模型
% \end{center}

% 主要内容：

% \begin{enumerate}
%   \item 数学规划模型的概念
%   \item 线性规划的求解方法
%   \item 非线性规划模型
%   \item 自来水运输问题
% \end{enumerate}

% \begin{center}
%     第七周\ 微分方程模型
% \end{center}

% 主要内容：

% \begin{enumerate}
%   \item Matlab求函数的导数、极限
%   \item Matlab求定积分、不定积分
%   \item Matlab求解微分方程
%   \item 药物在体内的分布模型
% \end{enumerate}

% \begin{center}
%     第八周\ 智能优化算法
% \end{center}

% 主要内容：

% \begin{enumerate}
%   \item 遗传算法
%   \item 免疫算法
%   \item 动态规划
%   \item 最短路径规划问题
%   \item 数学建模论文撰写规范
% \end{enumerate}

\section{建模案例}

\begin{center}
    双层玻璃窗保温原理的数学模型
\end{center}

\cepubblue{一. 问题提出}

    北方城镇楼房的窗户通常都是双层玻璃的，中间有一定空隙，双层玻璃比单层玻璃更保暖，那么到底比单层玻璃能减少多少热量损失呢？

\cepubblue{二. 问题分析}

    \begin{enumerate}
      \item 为了具有比较意义，双层玻璃与单层玻璃的厚度应该相同。
      \item 室内温度高，室外温度低，热量从室内往室外流失。
      \item 流失热量与窗户面积成正比，考虑总量没有意义，应该考虑单位面积流失的热量。
      \item 热传导需要考虑温度差、介质厚度，用变量把这些量表示出来。
    \end{enumerate}

    明确问题：
    \begin{enumerate}
      \item 根据热传导规律分别对单层玻璃、双层玻璃的热量流失规律建模，得到单位面积热量流失量，做出相对比较。
      \item 建立双层玻璃流失热量与空气层厚度的函数关系，合理地确定最优空气层厚度。
    \end{enumerate}

\cepubblue{三. 模型假设}

\begin{enumerate}
    \item 热量传递过程只有从室内到室外的单方向传导，没有对流辐射 (即不考虑热量流失到墙体等)。
    \item 室内温度与室外温度恒定不变，热传导处于稳定状态，即沿热传导方向，单位时间通过单位面积的热量是常数。
    \item  玻璃材质均匀，即导热系数 (也称热导率) 是常数。
  \end{enumerate}


\cepubblue{四. 建立模型}

热传导速率：
\[
    Q = -k \frac{\Delta T}{d}
\]

\begin{figure}[htpb]
    \centering
    \includegraphics[width=.4\linewidth]{双层玻璃}
    \caption{双层玻璃模型}
\end{figure}

双层玻璃模型中，热量经历3次传递，每次传递热量相等。
\[
    Q_2 = k_g \frac{T_1-T_a}{d} = k_a \frac{T_a-T_b}{l} = k_g \frac{T_b-T_2}{d}
\]

\begin{figure}[htpb]
    \centering
    \includegraphics[width=.4\linewidth]{单层玻璃}
    \caption{单层玻璃模型}
\end{figure}

单层玻璃模型热量流失为：
\[
    Q_1 = k_g \frac{T_1-T_2}{2d}
\]

\cepubblue{五. 模型求解}

双层玻璃模型为3个方程构成的方程组，用Matlab求解。
\[
    \begin{cases}
        Q_2 = k_g \frac{T_1-T_a}{d} \\
        Q_2 = k_a \frac{T_a-T_b}{l} \\
        Q_2 = k_g \frac{T_b-T_2}{d}
    \end{cases}
\]
Matlab代码如下：
\begin{lstlisting}
    syms kg ka T1 T2 l d Q2 Ta Tb; % 定义符号变量
    % 定义方程组
    eqns = [Q2 == kg*(T1-Ta)/d,
            Q2 == ka*(Ta-Tb)/l,
            Q2 == kg*(Tb-T2)/d];
    % 求解方程组
    S = solve(eqns,[Q2,Ta,Tb]);
    simplify(S.Q2)
\end{lstlisting}

计算结果：
\[
    Q_2=\frac{k_a k_g (T_1 - T_2)}{2 d k_a + lk_g }
\]

\cepubblue{六. 结果分析}

为了便于比较，将双层玻璃模型和单层玻璃模型流失热量方程形式一致化：
\[
    Q_1 = \frac{k_g(T_1-T_2)}{d(2+0)}
\]
\[
    Q_2 = \frac{k_g(T_1-T_2)}{2\left(2+\frac{l}{d}\frac{k_g}{k_a}\right)
    }
\]

双层玻璃和单层玻璃流失热量相比：
\[
    \frac{Q_2}{Q_1} = \frac{2}{2+s}
\]

其中，
\[
    s=\frac{l}{d}\frac{k_g}{k_a}
\]

查资料确定$k_g= 4 \sim 8 \times 10^{-3}$ \si{W/(m.K)}，$k_a= 2.5 \times 10^{-3}$ \si{W/(m.K)}，引入$h=l/d$(空气层与一层玻璃厚度之比)减少变量数，

\[
    r = \frac{Q_2}{Q_1} = \frac{2}{2+16h} = \frac{1}{1+8h}
\]

Matlab绘图展示：

\begin{figure}[htpb]
    \centering
    \includegraphics[width=.56\linewidth]{热量损失绘图}
\end{figure}

Matlab代码如下：

\begin{lstlisting}
r = @(h) 1 ./(1+8*h) % 定义匿名函数
h = 0:0.1:8;
plot(h,r(h),'-')
xlabel('$l/d$',Interpreter='latex'),ylabel('$Q_2/Q_1$',Interpreter='latex')
title('热量损失比Q2/Q1与l/d的关系')
\end{lstlisting}

比值$h$取多少最优呢？可以根据效果基本不再改进选择$h$，比如$\frac{\Delta r}{\Delta h}<0.01$：

\begin{lstlisting}
eqn = 1./(1+8*x);
h = double(solve(abs(diff(eqn)) == 0.01))
\end{lstlisting}

解得，$h=3.41$，此时$r=0.0354$，即可以减少热量损失96.46\%。

实际上，房间热量还可以通过墙壁、地面、顶棚等流失，双层玻璃功效会比模型结果差一些。


\end{sloppypar}
\end{document}